\ifdefined\tracebenchDocsCombinedAppendixMode
\def\tracebenchAppendixExtendedRelatedWorkStart{%
  \docTitle{Related Work}%
  \phantomsection\label{docstart:appendix_extended_related_work}%
  \section*{Related Work}%
}
\def\tracebenchAppendixExtendedRelatedWorkEnd{}
\else\ifdefined\tracebenchCombinedAppendixMode
\def\tracebenchAppendixExtendedRelatedWorkStart{%
  \section{Extended Related Work}%
  \label{app:extended-related-work}%
}
\def\tracebenchAppendixExtendedRelatedWorkEnd{}
\else
\documentclass[11pt]{article}
\usepackage{manual_docs_style}

\def\tracebenchAppendixExtendedRelatedWorkStart{%
  \begin{document}%
  \docTitle{Related Work}%
  \phantomsection\label{docstart:appendix_extended_related_work}%
  \section*{Related Work}%
}
\def\tracebenchAppendixExtendedRelatedWorkEnd{\end{document}}
\fi\fi
\newcommand{\relatedWorkMeta}[8]{%
  \paragraph{#1.}%
  \leavevmode\par\nobreak\vspace{0.05em}%
  \noindent{\scriptsize\setlength{\parskip}{0pt}%
    Citation: #2.
    First online: #3.
    Venue: #4.
    Dataset: #5.
    Availability: #6.
    Time-series length: #7.
    Task goal: #8.
    \\}%
}

\newcommand{\relatedWorkTextMeta}[3]{%
  \noindent{\scriptsize\setlength{\parskip}{0pt}%
    Text fields: #1.
    Text--series relation: #2.
    Text scope: #3.
    \par}%
  \vspace{0.2em}%
}



\tracebenchAppendixExtendedRelatedWorkStart

This appendix provides an analysis of dataset and benchmark papers that share similarity with \name{}.
It follows the convention mentioned in the \href{../../docs-template/tex/source/example.tex}{\code{docs-template/tex/source/example.tex}} file 
However with the additional addition of 
\begin{itemize}
  \item whether the data used to create the benchmark is new, reused, derived, or reframed from previous publications
  \item availability of the dataset with URL 
  \item time-series length, channel count, or source-dependent length information when applicable
  \item task goal
\end{itemize}
Additionally in each text of the time series it should be mention:
\begin{itemize}
  \item Only if it's the case paragraph with written "Reasons not to mention". 
  If this is the case nothing else is necessary.
  \item How was the dataset constructed 
  \item Weakness at the end (in respect to our paper)
\end{itemize}

\subsection*{Time-series and text benchmarks}


\relatedWorkMeta
  {\href{https://doi.org/10.1038/s41597-020-0495-6}{\emph{PTB-XL}}}
  {1419 (queried June 2, 2026)}{May 2020}{Nature Scientific Data}{New dataset}{PhysioNet, \url{https://physionet.org/content/ptb-xl/}}{10-second 12-lead ECG records}{Diagnostic classification}
It is a public clinical ECG dataset with real 12-lead electrocardiogram recordings and diagnostic annotations. It includes 21,837 10-second 12-lead ECG records from 18,885 patients.
A lead is the electrical activity of the heart measured from a certain angle.
A 12-lead electrocardiogram involves 10 electrodes from which you can get this.
It was created by a large collection of data in hosptials. 

\paragraph{Weakness}: Text wise it only includes metadata such sex, gender etc plus the ground truth target (i.e. SCP-ECG statement, which is the diseases (diagnosis) affecting the heart or a description)


\relatedWorkMeta
  {\href{https://proceedings.neurips.cc/paper_files/paper/2023/hash/d0b67349dd16b83b2cf6167fb4e2be50-Abstract-Datasets_and_Benchmarks.html}{\emph{ECG-QA}}}
  {23* (queried June 2, 2026)}{Jun 2023}{NeurIPS 2023 Datasets and Benchmarks Track}{New QA annotations/templates over existing ECG datasets}{Official repository, \url{https://github.com/Jwoo5/ecg-qa}}{10-second ECG records where source ECG datasets provide 10-second records; otherwise varies by source dataset}{Question answering}
This work reuses ECG signals and labels from existing ECG datasets, especially PTB-XL in its original version, and converts the clinical ECG labels into question-answer using templates.
Experts are not involved, they just convert ground truth into text with templates.

\paragraph{Weakness}: It inhereits all the weakness from the other work.




\relatedWorkMeta
  {\href{https://doi.org/10.1038/s41597-022-01899-x}{\emph{MIMIC-IV}}}
  {3305 (queried June 2, 2026)}{Jan 2023}{Scientific Data}{New clinical EHR dataset resource}{PhysioNet, credentialed access required}{Variable irregular clinical time series by patient stay and measurement type}{Clinical prediction, retrospective EHR analysis, and multimodal clinical modeling}
\relatedWorkTextMeta{Clinical notes, reports, labels, diagnoses, procedures, and structured clinical records}{Provides metadata, describes the patient state, and can provide external clinical context aligned to patient timelines}{Sample-level and patient-stay-level}
It is a large de-identified clinical electronic health record dataset with patient-stay measurements, structured records, and clinical documentation.
It is relevant because it combines irregular clinical time-series records with text and structured patient-stay context.


\relatedWorkMeta
  {\href{https://doi.org/10.13026/4nqg-sb35}{\emph{MIMIC-IV-ECG}}}
  {191 (queried June 8, 2026)}{Sep 2023}{PhysioNet dataset release}{New dataset resource/module linked to MIMIC-IV}{PhysioNet, \url{https://physionet.org/content/mimic-iv-ecg/1.0/}}{10-second 12-lead ECG records at 500 Hz}{Diagnostic interpretation}
\relatedWorkTextMeta{Diagnostic reports}{Describes the observation}{Sample-level}
It is a large hospital ECG resource with diagnostic reports matched to real clinical ECG waveforms.
It is relevant because it provides real-world medical signal-language data with paired waveform and report information.


\relatedWorkMeta
  {\href{https://arxiv.org/abs/2503.16858}{\emph{MTBench}}}
  {41 (queried June 2, 2026)}{Mar 2025}{arXiv preprint; ICLR 2026 OpenReview desk-rejected submission}{New benchmark dataset}{Official repository, \url{https://github.com/Graph-and-Geometric-Learning/MTBench}}{Varies by source dataset}{Temporal reasoning question answering}
\relatedWorkTextMeta{Question, context, and answer}{Queries the observation, provides external context, and is aligned to timestamps}{Sample-level}
They constructed the dataset as follows:
\begin{itemize}
    \item Finance side:
    \begin{itemize}
      \item Collected over 200,000 financial news article URLs from GlobeNews, MarketWatch, SeekingAlpha, Zacks,
      Invezz, Quartz, PennyStocks, and Benzinga.
      \item Covered May 2021--September 2023.
      \item Parsed each article's content, title, associated company/ticker names, and publication date.
      \item Curated 20,000 articles with balanced article lengths.
      \item Used GPT-4o to annotate content type, temporal effect range, and sentiment.
      \item Aligned each article with the corresponding market time series by publication timestamp and trend
      direction.
    \end{itemize}

    \item Weather side:
    \begin{itemize}
      \item Selected 50 U.S. airports from the GHCN-H hourly weather dataset.
      \item Covered 2003--2020.
      \item Used variables such as temperature, humidity, wind, visibility, pressure, and precipitation.
      \item Used NOAA's Storm Events Database for severe weather event type, location, casualties, and narrative
      descriptions.
      \item Aligned storm events to the nearest airport weather station within 50 km.
      \item Used LLMs to generate synthetic news-style articles for missing storm narratives.
    \end{itemize}

    \item Time series:
    \begin{itemize}
      \item Stock-market time series.
      \item Weather time series: temperature, humidity, wind, visibility, pressure, and precipitation.
    \end{itemize}
  \end{itemize}


\paragraph{Weakness}: Noisy data, not published (so we can avoid citing it), the questions are created by LLM.


\relatedWorkMeta
  {\href{https://arxiv.org/abs/2503.01875}{\emph{Time-MQA}}}
  {65 (queried June 2, 2026)}{Feb 2025}{ACL 2025 long paper}{New QA dataset, TSQA, derived from diverse time series}{Hugging Face dataset, \url{https://huggingface.co/datasets/Time-MQA/TSQA}}{Varies by source dataset}{Multi-task question answering, including forecasting, imputation, anomaly detection, classification, and open-ended reasoning}
\relatedWorkTextMeta{Question, context, and answer}{Queries the observation, provides task context, and is the target output}{Sample-level}

Existing time-series sample + task-specific natural-language template for classical tasks such as forecasting, imputation, anomaly detection, and classification $\rightarrow$ QA pair. 
For open-ended QA, existing time-series samples are paired with GPT-4o-generated questions and answers about properties such as trend, seasonality, volatility, anomalies, structural breaks, and summaries. 
For open ended QA Existing time-series sample + GPT4o questions about trend seasonaility etc. The reference answer is synthetic (GPT4o generated)

\paragraph{Weakness}: Questions are created by templates, or with GPT4o so ground truth is inherently biased and introduces circularity when testing.


\relatedWorkMeta
  {\href{https://openreview.net/forum?id=ULQt51DRug}{\emph{TRQA}}}
  {\texttt{NaN} (queried June 2, 2026)}{Sep 2025}{OpenReview; ICLR 2026 withdrawn submission}{New benchmark dataset}{\texttt{NaN}; no official dataset URL identified in the current notes}{Varies by source dataset/sample}{Time-series reasoning question answering, including anomaly detection and classification}
\paragraph{Reasons not to mention}: Not published anywhere and data not available.
It is a time-series reasoning question-answering benchmark that unifies multiple task families, including anomaly detection, classification, characterization, comparison, transformation, and temporal relationship reasoning.
It is relevant because it evaluates a broad set of temporal reasoning abilities in a QA format.


\relatedWorkMeta
  {\href{https://doi.org/10.1145/3715014.3722074}{\emph{SensorQA}}}
  {10 (queried June 8, 2026)}{Jan 2025}{ACM SenSys 2025}{Reused ExtraSensory sensor features plus new human-created QA annotations}{Official repository, \url{https://github.com/benjamin-reichman/SensorQA}}{Long-term sensor histories from a single day to multiple weeks}{Sensor question answering}
SensorQA builds on ExtraSensory sensor data, but its QA annotations are human-created rather than template-generated.
They asked  Amazon Mechanical Turk workers to generate QA from activity-history graphs such as this.
\begin{DocCode}
Monday
Sleeping:  #######
Walking:          ##
At work:            ########
Eating:                    ##
Sitting:              ######
\end{DocCode}
and provided answers derived from the sensor data.

\paragraph{Weakeness}: There is not much textual contex, it's about retrofitting an existing dataset into QA.


\relatedWorkMeta
  {\href{https://dl.acm.org/doi/10.1145/3450268.3453529}{\emph{DeepSQA}}}
  {12 (queried June 2, 2026)}{May 2021}{ACM/IEEE IoTDI 2021}{Generated OppQA benchmark over the reused OPPORTUNITY sensor dataset}{Official repository, \url{https://github.com/nesl/DeepSQA}}{Windowed OPPORTUNITY sensor sequences; exact length not specified in the current notes}{Sensor question answering}
\relatedWorkTextMeta{Question and answer}{Generated from labels, queries the observation, and is the target output}{Sample-level}
Dataset: Construction 
It introduces a sensory question-answering framework and the SQA-Gen pipeline for automatically generating QA datasets from labeled sensor data.
Using SQA-Gen, the authors create OppQA from the OPPORTUNITY human-activity sensor dataset.
The questions are generated from predefined functional templates over labeled activity windows, and the answers are computed from the corresponding ground-truth activity labels.
It is relevant because it studies natural-language question answering over wearable-sensor time series, although there are no textual descriptions.


\relatedWorkMeta
  {\href{https://arxiv.org/abs/2406.08627}{\emph{Time-MMD}}}
  {62 (queried June 2, 2026)}{Jun 2024}{NeurIPS 2024 Datasets and Benchmarks Track}{New multi-domain multimodal dataset}{Official repository, \url{https://github.com/AdityaLab/Time-MMD}}{Varies by source dataset/domain}{Forecasting}
\relatedWorkTextMeta{Context and answer}{Provides external context, is aligned to timestamps, and is the target output}{Sample-level}
It is a multi-domain multimodal time-series dataset that aligns numerical time series with textual information for multimodal time-series analysis.
It is relevant because it addresses the gap between numerical-only time-series modeling and real settings where text accompanies the series.


\relatedWorkMeta
  {\href{https://arxiv.org/abs/2509.05215}{\emph{BEDTime}}}
  {6 (queried June 2, 2026)}{Sep 2025}{arXiv preprint / OpenReview}{Benchmark unifying existing time-series description datasets}{\texttt{NaN}; no official dataset URL identified in the current notes}{Varies by source dataset; 46,843 is the number of time series/descriptions, not the per-series length}{Description recognition, description differentiation, and open-ended description generation}
\relatedWorkTextMeta{Description and generated text}{Describes the observation and is the target output}{Sample-level}
It is a benchmark for automatically describing time series through recognition, differentiation, and natural-language generation tasks.
It is relevant because it evaluates whether models can verbalize structural properties of time-series signals.


\relatedWorkMeta
  {\href{https://arxiv.org/abs/2509.20823}{\emph{CaTS-Bench}}}
  {4 (queried June 2, 2026)}{Sep 2025}{arXiv preprint; ACL 2026 Findings}{New benchmark dataset derived from 11 source datasets}{Official Hugging Face dataset page, \url{https://huggingface.co/datasets/mhfisher/CaTSBench}; public page visible, though the currently visible hosted data appears to expose a small subset}{Variable windows, 5--120 points in the visible hosted examples}{Captioning and time-series reasoning}
\relatedWorkTextMeta{Context, metadata, caption, question, and answer}{Provides metadata, describes the observation, queries the observation, and is the target output}{Sample-level}
It is a context-aware time-series captioning and reasoning benchmark.
It combines numeric time-series segments, contextual metadata, line-chart images, and natural-language captions across diverse real-world domains.

\paragraph{Issues}: (dataset not available)


\relatedWorkMeta
  {\href{https://github.com/y-kawagu/SUSHI}{\emph{SUSHI}}}
  {\texttt{NaN} (queried June 2, 2026)}{Sep 2024}{Zenodo dataset release}{New synthetic dataset}{Zenodo, \url{https://zenodo.org/records/13882998}; repository, \url{https://github.com/y-kawagu/SUSHI}}{2,048 points per signal}{Captioning and retrieval}
\relatedWorkTextMeta{Caption and description}{Describes the observation}{Sample-level}
It is a dataset of synthetic unichannel time-series signals paired with natural-language descriptions of signal characteristics.
It is relevant because it supports text-to-signal retrieval and signal captioning, directly connecting time-series patterns with language.


\relatedWorkMeta
  {\href{https://arxiv.org/abs/2411.06735}{\emph{Multi-Modal Forecaster}}}
  {26 (queried June 2, 2026)}{Nov 2024}{arXiv preprint}{New TimeText Corpus dataset}{Official repository, \url{https://github.com/Rose-STL-Lab/Multimodal_Forecasting}}{Variable timestamp-aligned sequences from climate and healthcare domains}{Joint multimodal forecasting}
\relatedWorkTextMeta{Context and answer}{Provides external context, is aligned to timestamps, and is the target output}{Sample-level}
It introduces the TimeText Corpus (TTC), a timestamp-aligned collection of textual and numerical time-series data for multimodal forecasting.
It is relevant because TTC represents the common setting where text is aligned to the time axis and may help forecast future numerical or textual values.


\relatedWorkMeta
  {\href{https://arxiv.org/abs/2506.09108}{\emph{SensorLM}}}
  {\texttt{NaN} (queried June 4, 2026)}{Jun 2025}{NeurIPS 2025}{New large-scale sensor-language dataset and model family}{Official repository, \url{https://github.com/Google-Health/consumer-health-research/tree/main/sensorlm}}{Wearable sensor sequences; over 59.7 million hours from more than 103,000 people}{Sensor-language modeling, activity analysis, healthcare, sensor captioning, and cross-modal retrieval}
\relatedWorkTextMeta{Generated hierarchical captions and task labels/descriptions}{Describes the observation and supports sensor-language alignment}{Sample-level}
It introduces SensorLM, a family of sensor-language foundation models trained with hierarchical captions that capture statistical, structural, and semantic information from wearable sensor data.
It is relevant because it studies how natural language can describe and retrieve wearable-sensor time series, but differs from \name{} because the language is primarily generated from observed sensor data rather than used to specify controlled simulator interventions.

\subsection*{How is \name{} different}
We have controlled hidden intervention that allows to have counterfactual reasoning.
No noise, characterisation of compleixirt.
Prompt goes to time sereis, arbistrary detailed allow for comleixty.

\subsection*{Data Analysis and ML Engineering Agents}

General data-analysis benchmarks often evaluate whether agents can load datasets, write analysis code, fit models, or answer questions about tables and plots.
These benchmarks are valuable but often the skill being measured is coding proficiency rather than data analysis, especially when task instructions are explicit or when the dataset is sufficiently large that generic machine-learning models perform well out of the box.

Moreover, many existing time-series datasets are poorly suited for evaluating whether agents can use a textual system description to guide direct inspection of time-series observations.
Specifically, in some settings, the time series are short enough to be serialized directly into the prompt, so the benchmark primarily tests single-pass reasoning rather than tool-mediated exploration.
In others, the accompanying text is timestamp-aligned, descriptive, or background context, rather than a reusable causal description of the data-generating system.
This makes it difficult to isolate whether an agent used the textual description to form hypotheses, inspect the relevant temporal evidence, and ground its final decision in the observed signals.
It thus remains unclear whether current tool-using agents can autonomously perform data analysis by combining textual descriptions with time-series observations.
TraceBench is a diagnostic framework designed to generate change-point and root-cause analysis tasks that can systematically assess agentic reasoning on time-series data.


\relatedWorkMeta
  {\href{https://arxiv.org/abs/2401.05507}{\emph{InfiAgent-DABench}}}
  {115 (queried June 2, 2026)}{Jan 2024}{ICML 2024}{New data-analysis benchmark/dataset, DAEval}{Project page, \url{https://infiagent.github.io/}; repository, \url{https://github.com/InfiAgent/InfiAgent}}{Not applicable}{Data analysis}
\relatedWorkTextMeta{Instruction and answer}{Not time-series-specific}{Task-level}
It evaluates tool-using agents on data-analysis tasks in an execution environment.
It is relevant because it measures whether agents can operate over data with tools rather than only answer static natural-language questions.


\relatedWorkMeta
  {\href{https://arxiv.org/abs/2407.06423}{\emph{InsightBench}}}
  {22 (queried June 2, 2026)}{Jul 2024}{arXiv preprint}{New business-analytics benchmark dataset}{Official repository, \url{https://github.com/ServiceNow/insight-bench}; project page, \url{https://insightbench.github.io/}}{Not applicable}{Insight generation}
\relatedWorkTextMeta{Instruction, context, and answer/insight summary}{Not time-series-specific}{Task-level}
It evaluates business analytics agents through multi-step insight-generation scenarios.
It is relevant because it studies whether agents can move from data inspection to useful analytical conclusions.


\relatedWorkMeta
  {\href{https://arxiv.org/abs/2409.07703}{\emph{DSBench}}}
  {90 (queried June 2, 2026)}{Sep 2024}{ICLR 2025}{New benchmark assembled from ModelOff and Kaggle tasks}{Official repository, \url{https://github.com/liqiangjing/dsbench}}{Not applicable}{Data-science workflow completion}
\relatedWorkTextMeta{Instruction and answer/solution}{Not time-series-specific}{Task-level}
It benchmarks data-science agents on tasks intended to approximate expert data-science workflows.
It is relevant because it evaluates end-to-end agent behavior across data preparation, modeling, and analysis.


\relatedWorkMeta
  {\href{https://arxiv.org/abs/2403.02528}{\emph{DACO}}}
  {9 (queried June 2, 2026)}{Mar 2024}{NeurIPS 2024 Datasets and Benchmarks Track}{New code-based data-analysis dataset built from tabular databases}{Official repository, \url{https://github.com/shirley-wu/daco}; project page, \url{https://shirley-wu.github.io/daco/}}{Not applicable}{Code-based data analysis}
\relatedWorkTextMeta{Instruction/query and answer}{Not time-series-specific}{Task-level}
It evaluates application-driven data analysis via code generation.
It is relevant because it treats code as the mechanism by which agents perform concrete analysis steps over data.


\relatedWorkMeta
  {\href{https://arxiv.org/abs/2402.17453}{\emph{DS-Agent}}}
  {114 (queried June 2, 2026)}{Feb 2024}{ICML 2024}{Benchmark dataset for automated data-science tasks}{Official repository, \url{https://github.com/guosyjlu/DS-Agent}}{Not applicable}{Automated data science}
\relatedWorkTextMeta{Instruction and answer/solution}{Not time-series-specific}{Task-level}
It studies automated data science by augmenting language models with case-based reasoning.
It is relevant because it explores how prior cases and structured workflows can guide agents through data-science tasks.


\relatedWorkMeta
  {\href{https://arxiv.org/abs/2410.07095}{\emph{MLE-bench}}}
  {241 (queried June 2, 2026)}{Oct 2024}{arXiv preprint}{Benchmark over 75 Kaggle ML engineering competitions}{Official repository, \url{https://github.com/openai/mle-bench}}{Not applicable}{ML engineering}
\relatedWorkTextMeta{Instruction and answer/submission}{Not time-series-specific}{Task-level}
It evaluates machine-learning agents on realistic ML engineering tasks.
It is relevant because it measures whether agents can execute practical model-development workflows rather than only produce isolated predictions.

\subsection*{Agentic coding benchmarks}


\relatedWorkMeta
  {\href{https://arxiv.org/abs/2605.07769}{\emph{Coding Agents Don't Know When to Act}}}
  {0 (queried June 2, 2026)}{May 2026}{arXiv preprint}{New FixedBench benchmark derived from human-verified coding tasks where no code change is required}{\texttt{NaN}; no official benchmark URL identified in the current notes}{Not applicable}{Bug-report action/inaction decision}
\relatedWorkTextMeta{Bug report and action decision}{Not time-series-specific}{Repository-level}
It evaluates coding agents on stale bug reports where the correct response is often to make no code change.
The benchmark is relevant because it tests whether an agent can resist an unsupported action rather than always producing an edit.

\begin{table}[ht]
\centering
\footnotesize
\setlength{\tabcolsep}{3pt}
\renewcommand{\arraystretch}{1.12}
\caption{Reported FixedBench run accounting for \emph{Coding Agents Don't Know When to Act}.
The SORCAR total is shown as 200--400 because it is 200 for the directly comparable setup, or 400 if the extra SORCAR prompt comparison is counted separately.}
\label{tab:coding_agents_act_fixedbench_runs}
\begin{tabular}{@{}>{\raggedright\arraybackslash}p{0.16\linewidth}>{\raggedright\arraybackslash}p{0.14\linewidth}>{\raggedright\arraybackslash}p{0.18\linewidth}>{\raggedright\arraybackslash}p{0.21\linewidth}>{\raggedright\arraybackslash}p{0.17\linewidth}@{}}
\hline
\textbf{Model} &
\textbf{Harness} &
\shortstack[l]{\textbf{Main}\\\textbf{FixedBench}\\\textbf{runs}} &
\shortstack[l]{\textbf{Extra Table 1}\\\textbf{ablations?}} &
\shortstack[l]{\textbf{Total counted}\\\textbf{runs}} \\
\hline
Sonnet-4.6 & Claude Code & included & Yes & 1,700 \\
GPT-5.4 mini & Codex & included & Yes & 1,700 \\
GPT-5.3 Codex & Codex & 200 & No & 200 \\
Gemini-3 Pro & Gemini-CLI & 200 & No & 200 \\
Qwen3.5-122B & Qwen-Code & 200 & No & 200 \\
GPT-5.3 Codex & SORCAR & 200 & No / limited comparison & 200--400 \\
\hline
\end{tabular}
\end{table}


\relatedWorkMeta
  {\href{https://arxiv.org/abs/2602.11988}{\emph{Evaluating AGENTS.md}}}
  {7 (queried June 2, 2026)}{Feb 2026}{arXiv preprint}{Evaluation set reusing SWE-bench tasks plus repository issues/context files}{\texttt{NaN}; no official dataset URL identified in the current notes}{Not applicable}{Repository-instruction following}
\relatedWorkTextMeta{Repository instruction and task}{Not time-series-specific}{Repository-level}
It evaluates whether repository-level instruction files improve coding-agent behavior.
The paper is relevant because it studies how additional context changes agent performance, tool use, and decision quality.


\relatedWorkMeta
  {\href{https://arxiv.org/abs/2603.04177}{\emph{CodeTaste}}}
  {0 (queried June 2, 2026)}{Mar 2026}{arXiv preprint}{New benchmark mined from large-scale multi-file open-source refactoring changes}{Official repository, \url{https://github.com/logic-star-ai/refactoring-benchmark}}{Not applicable}{Code refactoring}
\relatedWorkTextMeta{Refactoring task specification and answer/patch}{Not time-series-specific}{Repository-level}
It evaluates whether language models can perform human-level code refactorings in realistic codebases, including under weaker or less prescriptive task specifications.
The benchmark is relevant because it probes whether an agent can move beyond following detailed instructions and produce a useful reusable change.

\subsection*{Dataset Examples}

\phantomsection\label{app:related-work-dataset-example-mtbench}
\paragraph{\emph{MTBench}.}
\textbf{Q:}
\begin{DocCode}
Given the reported intense weather conditions on April 13th, including hail and blizzard scenarios, which of the following options about the expected temperature patterns on April 14th is most accurate?

A) Following the blizzard conditions, a sudden increase in temperature is expected, leading to a rapid thaw and the potential for flooding due to melting snow.  
B) After the hailstorm, temperatures are likely to stabilize around freezing with limited fluctuation, preventing any significant thawing.  
C) The significant drop in temperature after April 13th indicates the likelihood of a second storm system moving in from the northwest, which could bring additional snowfall.  
D) The temperature is expected to drop due to the high winds and blizzard conditions experienced on the previous day, leading to a continued sub-freezing environment.

Context:
The following events were reported: Hail. These occurred near station USW00003017, approximately 27.8164 km away, between 2018-04-13 00:45:00 and 2018-04-13 00:45:00. The events included records with hailstones measuring 0.75 inches.Thankfully, no injuries or fatalities were reported. No significant property or crop damage was reported. Episode Narrative: A powerful storm system moved over northeast Colorado on the 13th, then slowly exited into the Central Plains on the 14th. The storm produced blizzard conditions over the far northeastern plains of Colorado, with high winds elsewhere. Snow and blowing snow produced extensive drifting snow up to 7 feet deep, along with near zero visibility. The dangerous conditions made the roads impassable and stranded motorists over northeast Lincoln, Logan, Phillips, Sedgwick, and Washington counties. Consequently, I-76 east of Sterling, I-70 east of Limon and nearly all other roads and highways across Northeast and East-Central Colorado were closed due to hazardous driving conditions. In total, more than a dozen semi-trucks were blown over. Eight semi trucks west of Akron, along Colorado Highway 34 were blown over. Blizzard conditions made it nearly impossible to get tow trucks to accidents and rescue stranded motorists. The dangerous conditions also caused widespread power outages and strayed cattle. Some livestock was lost, mainly calves.||Numerous trees were snapped or uprooted in the Hugo and Limon areas. A large tree, 2.5 ft in diameter, was knocked over in Hugo. Peak wind gusts included:  78 mph at  Akron-Washington County Airport, 76 mph at Limon Municipal Airport, 72 mph in Hugo and Sterling, 71 mph near Bennett and Genoa, 65 mph at Briggsdale, 62 mph, 7 miles northwest of Boyero and Denver International Airport; 60 mph, 8 miles south of Holyoke and near Wiggins; 59 mph near Keenesburg, 12 miles east-northeast of Kiowa, 7 miles south of Shamrock and Woodlin School; 58 mph at Agate, Front Range Airport in Watkins, Kassler and 11 miles west-northwest of Westplains. Snowfall totals ranged from 4 to 10 inches. Event Narrative:
\end{DocCode}

\textbf{A:}
\begin{DocCode}
D
\end{DocCode}

\phantomsection\label{app:related-work-dataset-example-time-mqa}
\paragraph{\emph{Time-MQA}.}
\textbf{Q:}
\begin{DocCode}
Identify any anomalies within the data set [8.0, 8.0, 5.0, 13.0, 6.0, 10.0, 20.0, 24.0, 14.0, 2.0, 21.0, 16.0, 21.0, 6.0, 11.0, 14.0, 35.0, 23.0, 0.0, 0.0, 0.0, 10.0, 33.0, 5.0].
\end{DocCode}

\textbf{A:}
\begin{DocCode}
The sequence '0.0, 0.0, 0.0' is an anomaly considering the surrounding higher values in the dataset, indicating an unusual drop in the series.
\end{DocCode}

\phantomsection\label{app:related-work-dataset-example-sensorqa}
\paragraph{\emph{SensorQA}.}
\textbf{Q:}
\begin{DocCode}
What did I do outside?
\end{DocCode}

\textbf{A:}
\begin{DocCode}
You walked.
\end{DocCode}

\phantomsection\label{app:related-work-dataset-example-time-mmd}
\paragraph{\emph{Time-MMD}.}
\textbf{Q:}
\begin{DocCode}
Use the agriculture numerical time series and the paired market report to produce the short-term textual forecast.

Context:
Whole broiler/fryer prices are steady, and retail demand is light to moderate. Supplies of all sizes are light to moderate.
\end{DocCode}

\textbf{A:}
\begin{DocCode}
Based on the current situation, it is likely that the Monthly Retail broiler Composite will remain stable or experience a slight increase in the next 1-3 months.
\end{DocCode}

\phantomsection\label{app:related-work-dataset-example-cats-bench}
\paragraph{\emph{CaTS-Bench}.}
\textbf{Q:}
\begin{DocCode}
Here is a time series:
37.00,37.00,37.00,37.00,37.00,40.57,40.57,40.57,40.57,40.57,40.57


What caption best describes to this time series?
(A) From May 1st to July 26th, 2024, the daily COVID-19 deaths in China show a fluctuating pattern, with values generally ranging between 0.29 and 2.86. There are periods of relative stability, such as the initial days of May with a consistent 0.86, interspersed with occasional spikes to 2.86, and dips to 0.29 towards the end of July. Compared to the general daily death statistics for China, where the mean is 73.0 and the maximum reaches 6812.0, this specific time series indicates a period of significantly lower daily deaths, suggesting a substantial improvement in the COVID-19 situation during this timeframe.

(B) From October 24, 2023, to November 3, 2023, the daily COVID-19 cases in Luxembourg show a relatively stable pattern, beginning at 37.3 cases and rising to 40.57 cases by October 29, 2023, where it remains for the rest of the period. Compared to the country's general statistics, where the mean is 236, the daily cases during this period are significantly lower, suggesting a period of reduced viral transmission. This trend does not follow any expected seasonal patterns, as COVID-19 case numbers are known to fluctuate unpredictably.

(C) From October 24, 2023, to November 3, 2023, the daily COVID-19 cases in Luxembourg show a relatively stable pattern, beginning at 37 cases and rising to 40.57 cases by October 29, 2023, where it remains for the rest of the period. Compared to the country's general statistics, where the mean is 236.0, the daily cases during this period are significantly lower, suggesting a period of reduced viral transmission. This trend does not follow any expected seasonal patterns, as COVID-19 case numbers are known to fluctuate unpredictably.

(D) From October 24, 2023, to November 3, 2023, the daily COVID-19 cases in Luxembourg show a relatively unstable pattern, beginning at 37 cases and decreasing to 40.57 cases by October 29, 2023, where it remains for the rest of the period. Compared to the country's general statistics, where the mean is 236.0, the daily cases during this period are significantly lower, suggesting a period of reduced viral transmission. This trend does follow expected seasonal patterns, as COVID-19 case numbers are known to fluctuate unpredictably.


You must respond only with valid JSON, and no extra text or markdown.
The JSON schema is:
{
  "answer": <string>
}
<string> must be an answer string containing only A, B, C, or D.
Ensure your output parses as JSON with exactly one top-level object containing the answer field.
\end{DocCode}

\textbf{A:}
\begin{DocCode}
C
\end{DocCode}

\phantomsection\label{app:related-work-dataset-example-sushi}
\paragraph{\emph{SUSHI}.}
\textbf{Q:}
\begin{DocCode}
Caption this unichannel synthetic time-series signal.
\end{DocCode}

\textbf{A:}
\begin{DocCode}
The trend remains constant, with no change from a specific value.
\end{DocCode}

\phantomsection\label{app:related-work-dataset-example-multi-modal-forecaster}
\paragraph{\emph{Multi-Modal Forecaster}.}
\textbf{Q:}
\begin{DocCode}
Given the first seven daily climate rows and their aligned forecast texts, forecast the next day numerical weather values.
\end{DocCode}

\textbf{A:}
\begin{DocCode}
{"date":"2014-01-08","temp":"22.2","precip":"0.0","humidity":"40.4","windspeed":"7.8"}
\end{DocCode}


\tracebenchAppendixExtendedRelatedWorkEnd
